% sec:sr-lorentz — Lorentz Transformations, Boosts, and Observer Tetrads
\section{Lorentz Transformations, Boosts, and Observer Tetrads}
\label{sec:sr-lorentz}

Two physicists in relative motion measure the same pair of events and compute
the spacetime interval.  Both must get the same answer---that is the content of
the invariance principle from \cref{sec:sr-minkowski}.  But they use different
coordinates, so their individual numbers for $\Delta t$ and $\Delta x$ differ.
Which coordinate changes are compatible with this requirement?  The answer is
the Lorentz group, and its most physically important elements---boosts---are
the transformations that connect observers in relative motion.  The tetrad
construction at the end of this section packages an observer's local frame into
four vectors that generalise directly to curved spacetime, where they define the
ray tracer's camera.

\paragraph{The Lorentz group.}

A \emph{Lorentz transformation} is a linear map $x'^\mu = \Lambda^\mu{}_\nu\,
x^\nu$ that preserves the Minkowski interval.  The preservation condition is
that for any pair of events,
%
\begin{equation}
  \eta_{\mu\nu}\,\d x'^\mu\,\d x'^\nu
    = \eta_{\mu\nu}\,\d x^\mu\,\d x^\nu \,.
  \label{eq:interval-invariance}
\end{equation}
%
Substituting $\d x'^\mu = \Lambda^\mu{}_\alpha\,\d x^\alpha$ and requiring this
to hold for arbitrary $\d x^\alpha$, we obtain the defining condition
%
\begin{equation}
  \eta_{\mu\nu}\,\Lambda^\mu{}_\alpha\,\Lambda^\nu{}_\beta
    = \eta_{\alpha\beta} \,.
  \label{eq:lorentz-condition}
\end{equation}
%
Read this as a constraint on the matrix $\Lambda$: the Minkowski metric is
``sandwiched'' between two copies of $\Lambda$ and must reproduce itself.  This
is the hyperbolic analogue of the orthogonality condition $R^T R = I$ in
Euclidean geometry---there, the identity matrix plays the role that $\eta$ plays
here.  In matrix notation, \cref{eq:lorentz-condition} reads
%
\begin{equation}
  \Lambda^{T}\!\eta\,\Lambda = \eta \,.
  \label{eq:lorentz-matrix}
\end{equation}
%
Any $4 \times 4$ real matrix satisfying this condition belongs to the
\emph{Lorentz group} $O(1,3)$.
\histnote{The group is named after Hendrik Lorentz, who derived the
transformations in 1904.  Einstein's insight was that they express a physical
symmetry, not merely a mathematical convenience.}

Taking the determinant of both sides of \cref{eq:lorentz-matrix} gives
$(\det\Lambda)^2 = 1$, so $\det\Lambda = \pm 1$.  Setting $\alpha = \beta = 0$
in \cref{eq:lorentz-condition} gives
$-(\Lambda^0{}_0)^2 + \sum_i (\Lambda^i{}_0)^2 = -1$; since the spatial sum
$\sum_i (\Lambda^i{}_0)^2$ is non-negative, it follows that
$(\Lambda^0{}_0)^2 \geq 1$.  The \emph{proper orthochronous} Lorentz group
$SO^+(1,3)$ consists of those elements with $\det\Lambda = +1$ and
$\Lambda^0{}_0 \geq 1$---these are the physically realisable transformations
that preserve spatial orientation and the direction of time.  We restrict
attention to this subgroup.

\paragraph{Boosts.}

The simplest proper orthochronous Lorentz transformations that mix space and
time are \emph{boosts}.  Consider two observers: $S$, at rest, and $S'$, moving
with velocity $v$ in the $x$-direction relative to $S$ (with $c = 1$, so
$v \in [0,1)$).  The transformation connecting their coordinates must satisfy
\cref{eq:lorentz-condition}, leave $y$ and $z$ unchanged, and reduce to the
identity at $v = 0$.  The unique solution is
%
\begin{equation}
  \Lambda^\mu{}_\nu
    = \begin{pmatrix}
        \gamma   & -v\gamma & 0 & 0 \\
        -v\gamma & \gamma   & 0 & 0 \\
        0        & 0        & 1 & 0 \\
        0        & 0        & 0 & 1
      \end{pmatrix} ,
  \label{eq:boost-matrix}
\end{equation}
%
where the \emph{Lorentz factor} is
%
\begin{equation}
  \gamma = \frac{1}{\sqrt{1 - v^2}} \,.
  \label{eq:lorentz-factor}
\end{equation}
%
This is the passive transformation giving $S'$ coordinates in terms of $S$
coordinates, with $S'$ moving at velocity $+v$ along $x$.  At $v = 0$ the
matrix is the identity; as $v \to 1$ the off-diagonal entries diverge,
reflecting the impossibility of boosting to the speed of light.
\warnnote{The signs in \cref{eq:boost-matrix} depend on the signature
convention.  In $(+{-}{-}{-})$ sources, the off-diagonal entries have
the opposite sign.}

It is worth verifying \cref{eq:lorentz-condition} explicitly.  The $(0,0)$
component gives $-\gamma^2 + v^2\gamma^2 = -\gamma^2(1 - v^2) = -1$, which
holds by definition of $\gamma$.  The $(0,1)$ component gives $v\gamma^2 -
v\gamma^2 = 0$, confirming the off-diagonal entries of $\eta$ vanish.  The
spatial block reduces to the identity since $\Lambda$ acts trivially on $y$ and
$z$.

\paragraph{Rapidity.}

An alternative parametrisation replaces $v$ with the \emph{rapidity} $\phi$,
defined by $v = \tanh\phi$ (equivalently, $\gamma = \cosh\phi$ and $v\gamma =
\sinh\phi$).  The boost matrix becomes
%
\begin{equation}
  \Lambda^\mu{}_\nu
    = \begin{pmatrix}
        \cosh\phi  & -\sinh\phi & 0 & 0 \\
        -\sinh\phi & \cosh\phi  & 0 & 0 \\
        0          & 0          & 1 & 0 \\
        0          & 0          & 0 & 1
      \end{pmatrix} .
  \label{eq:boost-rapidity}
\end{equation}
%
The payoff is that rapidities add: two successive collinear boosts with
rapidities $\phi_1$ and $\phi_2$ compose to a single boost with rapidity
$\phi_1 + \phi_2$.  This follows from the hyperbolic addition formula:
$\cosh(\phi_1 + \phi_2) = \cosh\phi_1\cosh\phi_2 + \sinh\phi_1\sinh\phi_2$,
which is exactly the $(0,0)$ entry of the matrix product
$\Lambda(\phi_1)\,\Lambda(\phi_2)$.  Contrast this with the relativistic
velocity-addition formula $w = (v_1 + v_2)/(1 + v_1 v_2)$---velocities compose
nonlinearly, but their corresponding rapidities simply sum.  For non-collinear
boosts, the composition also produces a spatial rotation (the Wigner rotation).
\physnote{Rapidity is the hyperbolic analogue of the rotation angle.  Rotations
compose by adding angles; boosts compose by adding rapidities.}

\paragraph{Observer tetrads.}

A Lorentz transformation connects two coordinate systems.  A \emph{tetrad}
goes further: it is an observer's personal ruler-and-clock kit planted at a
single point in spacetime.  One vector says ``this is my future direction,'' and
three mutually perpendicular vectors say ``these are my spatial axes.''  The
formal requirement is that these four vectors reproduce the Minkowski metric
when measured against the background.

An \emph{orthonormal tetrad} is a set of four vectors
$e^\mu{}_{(a)}$, with $a = 0, 1, 2, 3$, satisfying
%
\begin{equation}
  \eta_{\mu\nu}\, e^\mu{}_{(a)}\, e^\nu{}_{(b)}
    = \eta_{ab} \,.
  \label{eq:tetrad-orthonormality}
\end{equation}
%
The parenthesised index $(a)$ is a \emph{frame index}: it labels which tetrad
vector we are referring to, not a spacetime direction.  The notation
$e^\mu{}_{(a)}$ carries one contravariant spacetime index $\mu$ (giving the
vector's components in coordinates) and one covariant frame index $(a)$
(labelling which basis vector); the inverse tetrad $e^{(a)}{}_\mu$ reverses
these roles.  The orthonormality condition says that the timelike vector
$e^\mu{}_{(0)}$ has norm $-1$ and is orthogonal to the three spacelike vectors
$e^\mu{}_{(i)}$, each of norm $+1$.
\warnnote{Tetrad indices $(a), (b), \ldots$ live in a \emph{local} Minkowski
frame.  They are not raised or lowered with the spacetime metric---they use
$\eta_{ab}$ directly.}

For an observer at rest in the standard coordinates, the tetrad is trivial:
%
\begin{equation}
  \begin{aligned}
    e^\mu{}_{(0)} &= (1,0,0,0), \quad &
    e^\mu{}_{(1)} &= (0,1,0,0), \\
    e^\mu{}_{(2)} &= (0,0,1,0), \quad &
    e^\mu{}_{(3)} &= (0,0,0,1) \,.
  \end{aligned}
  \label{eq:tetrad-rest}
\end{equation}
%
This is the identity matrix---unsurprising, since the observer's frame
coincides with the coordinate frame.  For an observer moving with velocity
$v$ in the $x$-direction, we apply the inverse boost (replacing $v \to -v$ in
\cref{eq:boost-matrix}) to the rest-frame tetrad vectors, transforming them
from the moving observer's frame into $S$ coordinates:
%
\begin{equation}
  e^\mu{}_{(0)} = (\gamma,\, v\gamma,\, 0,\, 0), \quad
  e^\mu{}_{(1)} = (v\gamma,\, \gamma,\, 0,\, 0) \,.
  \label{eq:tetrad-boosted}
\end{equation}
%
The vectors $e^\mu{}_{(2)}$ and $e^\mu{}_{(3)}$ are unchanged.  The timelike
tetrad vector $e^\mu{}_{(0)} = \fourvel$ is the observer's
\emph{four-velocity}---this is not a coincidence but the formal definition of
``the direction the observer moves through spacetime.''  The tetrad is a
coordinate-free encoding of the observer's state of motion: it captures both
their velocity (via $e^\mu{}_{(0)}$) and their spatial orientation (via
$e^\mu{}_{(i)}$), which spans the observer's simultaneity surface and defines
the directions they call ``right,'' ``up,'' and ``forward.''

\paragraph{A worked example: camera frame at $v = 0.6$.}

Consider an observer (a camera, in our context) moving at $v = 0.6$ in the
$x$-direction.  The Lorentz factor is $\gamma = 1/\sqrt{1 - 0.36} = 1.25$.
The timelike tetrad vector is
$e^\mu{}_{(0)} = (1.25,\, 0.75,\, 0,\, 0)$,
and we can verify its norm:
$\eta_{\mu\nu}\,e^\mu{}_{(0)}\,e^\nu{}_{(0)} = -(1.25)^2 + (0.75)^2 =
-1.5625 + 0.5625 = -1$.  The spatial vector
$e^\mu{}_{(1)} = (0.75,\, 1.25,\, 0,\, 0)$ has norm
$-(0.75)^2 + (1.25)^2 = -0.5625 + 1.5625 = +1$, and the inner product
$\eta_{\mu\nu}\,e^\mu{}_{(0)}\,e^\nu{}_{(1)} = -(1.25)(0.75) + (0.75)(1.25)
= 0$.  The tetrad is indeed orthonormal.

\begin{figure}[htbp]
  \centering
  % boosted-tetrad.tex — Boosted tetrad vectors in 1+1 Minkowski space
% Inputable TikZ fragment for fig:boosted-tetrad
% Requires: tikz with calc, arrows.meta (loaded by ehbook.cls)
%
% Boost v = 0.6 → γ = 1.25
% e_{(0)} = (γ, γv) = (1.25, 0.75)   — timelike, tilted from t-axis
% e_{(1)} = (γv, γ) = (0.75, 1.25)   — spacelike, tilted from x-axis
% Tilt angle: arctan(0.6) ≈ 31°
%
\begin{tikzpicture}[
    >={Stealth[length=5pt]},
    font=\footnotesize,
  ]

  \def\axlen{2.8}     % axis half-length
  \def\cone{2.4}      % light-cone extent along each ray
  \def\vecscale{1.4}  % scale factor for tetrad vectors

  % ── Light-cone shading ────────────────────────────────────────────
  \fill[EHMarginGray, fill opacity=0.05]
    (0,0) -- (-\cone,\cone) -- (\cone,\cone) -- cycle;
  \fill[EHMarginGray, fill opacity=0.05]
    (0,0) -- (-\cone,-\cone) -- (\cone,-\cone) -- cycle;

  % ── Light-cone edges ──────────────────────────────────────────────
  \draw[EHMarginGray, thin] (-\cone,\cone)  -- (0,0) -- (\cone,\cone);
  \draw[EHMarginGray, thin] (-\cone,-\cone) -- (0,0) -- (\cone,-\cone);

  % ── Lab frame axes (S) ───────────────────────────────────────────
  \draw[->, thick, EHMarginGray] (-\axlen,0) -- (\axlen,0)
    node[right, font=\small] {$x$};
  \draw[->, thick, EHMarginGray] (0,-\axlen) -- (0,\axlen)
    node[above, font=\small] {$t$};

  % ── Boosted axes (dashed, to show tilted coordinate grid) ────────
  % t'-axis along e_{(0)} direction: slope = γ/(γv) = 1/v = 1/0.6
  \draw[EHMarginGray, thin, densely dashed]
    (0,0) -- ({0.6*\cone},\cone)
    node[right, font=\scriptsize, text=EHMarginGray] {$t'$};
  % x'-axis along e_{(1)} direction: slope = γv/γ = v = 0.6
  \draw[EHMarginGray, thin, densely dashed]
    (0,0) -- (\cone,{0.6*\cone})
    node[above, font=\scriptsize, text=EHMarginGray] {$x'$};

  % ── Timelike tetrad vector e_{(0)} (gold) ────────────────────────
  % e_{(0)} = (γ, γv) = (1.25, 0.75), plotted as (x,t) = (0.75, 1.25)
  \draw[->, very thick, EHAccretionGold]
    (0,0) -- ({0.75*\vecscale},{1.25*\vecscale})
    node[left=2pt, text=EHAccretionGold!85!black]
      {$e^\mu{}_{(0)}$};

  % ── Spacelike tetrad vector e_{(1)} (blue) ───────────────────────
  % e_{(1)} = (γv, γ) = (0.75, 1.25), plotted as (x,t) = (1.25, 0.75)
  \draw[->, very thick, EHJetBlue]
    (0,0) -- ({1.25*\vecscale},{0.75*\vecscale})
    node[right=2pt, text=EHJetBlue!85!black]
      {$e^\mu{}_{(1)}$};

\end{tikzpicture}

  \caption[Boosted tetrad vectors]{Tetrad vectors for an observer moving at $v = 0.6$ in the
    $x$-direction.  The timelike vector $e^\mu{}_{(0)}$ (gold) is the
    observer's four-velocity, tilted toward the $x$-axis by the boost.  The
    spatial vector $e^\mu{}_{(1)}$ (blue) is the observer's $x'$-direction,
    tilted symmetrically toward the $t$-axis.  Together they satisfy the
    orthonormality condition
    $\eta_{\mu\nu}\,e^\mu{}_{(a)}\,e^\nu{}_{(b)} = \eta_{ab}$.  In curved
    spacetime, the same construction defines the camera's local reference
    frame at any point.}
  \label{fig:boosted-tetrad}
\end{figure}

\paragraph{From flat to curved.}
\coderef{Spacetime}{Minkowski}

In the codebase, the Minkowski metric is an instance of the \texttt{Spacetime}
typeclass, with \texttt{metric4} returning \texttt{sym4Diag (-1) 1 1 1} and all
Christoffel symbols set to zero.  The same typeclass provides Schwarzschild and
Kerr metrics, and the camera's tetrad construction works identically for all of
them---only the orthonormality condition changes from $\eta_{\mu\nu}$ to
$\metric$.
\implnote{The \texttt{Spacetime} typeclass abstracts over the specific metric,
so the tetrad and ray-tracing code is metric-agnostic.}

In curved spacetime, the tetrad condition generalises to
%
\begin{equation}
  g_{\mu\nu}\,e^\mu{}_{(a)}\,e^\nu{}_{(b)} = \eta_{ab} \,.
  \label{eq:tetrad-curved}
\end{equation}
%
The metric on the left is the curved-spacetime metric, but the right-hand side
remains the flat Minkowski matrix.  The tetrad defines a locally Minkowski frame
at each point, regardless of the background curvature.  The camera model in
\cref{ch:camera-rendering} uses exactly this construction to define the
observer's ``sky''---the directions from which photons arrive---at any point in
any spacetime.
\xrefnote{The tetrad reappears in the camera model of
\cref{ch:camera-rendering}, where it defines the mapping from pixel
coordinates to initial photon momenta.}
